\documentclass[12pt,a4paper]{ctexart}

\usepackage[margin=2.5cm]{geometry}
\usepackage[colorlinks=true,linkcolor=black,urlcolor=blue]{hyperref}
\usepackage{enumitem}
\usepackage{longtable}

\title{基于局域网常开工控机与 Tailscale 的远程访问主机操作手册}
\author{大纲}
\date{\today}

\begin{document}

\maketitle
\tableofcontents
\clearpage



\section{目标拓扑}
\begin{itemize}[leftmargin=2em]
    \item 远程客户端：外部网络中的电脑或手机。
    \item 工控机：局域网内常开，安装 Tailscale，作为跳板或管理入口。
    \item 目标主机：真正需要远程控制的工作主机，需要通过网线连接网络。
    \item 路由器/交换机：提供局域网连接。
\end{itemize}



\section{步骤 1：部署工控机基础环境}

\subsection{操作内容}
\begin{itemize}[leftmargin=2em]
    \item 安装debian12系统，安装过程中需要选择安装ssh服务和标准系统工具。
    \begin{verbatim}
        进入root权限
        su -
        更新软件源、升级系统
        apt update
        apt upgrade -y
        安装SSH服务（如果没有安装的话）
        apt install -y openssh-server
        启动SSH服务
        systemctl start ssh
        设置SSH服务开机自启
        systemctl enable ssh
        检查SSH服务状态
        systemctl status ssh
        查看本机ip
        ip a

    \end{verbatim}
    \item 安装tailscale服务。
    \begin{verbatim}
        测试能否直接走tailscale官方服务
        curl -I --max-time 10 https://tailscale.com
        curl -I --max-time 10 https://login.tailscale.com
        curl -I --max-time 10 https://controlplane.tailscale.com
        如果上面都返回200或302，则表示能直接走官方服务
        安装tailscale
        curl -fsSL https://tailscale.com/install.sh | sh
        启动tailscale服务
        systemctl start tailscaled
        设置tailscale服务开机自启
        systemctl enable tailscaled
        在tailscale后台生成一个新的authkey，打开
        https://login.tailscale.com/admin/authkeys
        用auth key登录debian
        export TS_AUTH_KEY='生成的KEY'
        tailscale up --reset --auth-key="$TS_AUTH_KEY" --hostname=起一个名字（如debian） --ssh
        unset TS_AUTH_KEY
        查看tailscale是否登录成功
        tailscale status
        tailscale ip -4
        tailscaale netcheck
        正常时会看到
        tailscale status显示这台机器在线
        tailscale ip -4显示一个100开头的IP地址
        tailscale netcheck返回连通性信息
        确保开机自启动
        systemctl is-enabled tailscaled
        关闭后台的key expiry，实现长期在线

    \end{verbatim}
    \item 安装唤醒工具
    \begin{verbatim}
        安装wakeonlan工具
        apt install -y wakeonlan etherwake
        确认debian的局域网网段和接口
        ip a
        ip link
        根据他的ip地址确认广播地址，如192.168.10.109/24的广播地址是192.168.10.255
        准备目标主机的真实MAC地址，在目标主机执行ipconfig /all

    \end{verbatim}
\end{itemize}



\section{步骤 2：配置目标主机的远程访问能力}

\subsection{操作内容}
\begin{itemize}[leftmargin=2em]
    \item BOIS里开启远程唤醒功能（如果支持）。
    常见名字包括 Wake-on-LAN、Power on by PCI-E、Power on by LAN 等。
    \item windows设置
    \begin{verbatim}
    设备管理器-网络适配器-有线网卡-属性-高级
    wake on magic packet 设置为 enabled
    wake on pattern match 设置为 enabled
    shutdown wake on lan 设置为 enabled
    电源管理里勾选允许此设备唤醒计算机，只允许魔术包唤醒计算机
    关闭Windows快速启动
    控制面板-电源选项-选择电源按钮的功能-更改当前不可用的设置-取消勾选启用快速启动

    \end{verbatim}
    \item 安装tailscale客户端   
    \begin{verbatim}
    直接通过网页搜索软件安装即可，并登录同一账号，确认能看到工控机在线
    在debian主机上发送唤醒包测试
    wakeonlan -i 广播地址（如192.168.10.255） 目标主机MAC地址
    完整例子 wakeonlan -i 192.168.10.255 00:11:22:33:44:55

    \end{verbatim}
\end{itemize}





\end{document}
